\documentclass{article}
\usepackage{amsmath,amssymb,amsthm}
\usepackage{booktabs}
\usepackage{hyperref}
\usepackage[margin=1in]{geometry}
\usepackage{algorithm}
\usepackage{algpseudocode}

\newcommand{\hex}[1]{\texttt{\$#1}}
\newcommand{\Beff}{B_{\mathrm{eff}}}
\newcommand{\PASS}{\textsc{pass}}
\newcommand{\FAIL}{\textsc{fail}}

\newtheorem{definition}{Definition}
\newtheorem{proposition}{Proposition}

\title{Simulation-Trained Weighted Finite-State Transducers\\
for Replicator Viability Estimation}
\author{}
\date{March 2026}

\begin{document}
\maketitle

\begin{abstract}
We describe a weighted finite-state transducer (WFST) framework for
estimating the probability that a random byte stream encodes a viable
self-replicating program in a 6502 cellular automaton. The WFST
encodes structural constraints (opcode pattern, address alignment,
branch offset) as hard \FAIL\ verdicts, and learns soft transition
weights for ``slide'' positions from simulation. We derive an
AND-gate blame assignment for the EM M-step, validate with importance
sampling, and report $\Beff \approx 61$ bits for the minimal
replicator family. Active learning concentrates simulation budget
on uncertain transitions.
\end{abstract}

\section{The Scoring Problem}

Given a random byte stream of length $L$, what is the probability it
encodes a viable self-replicating 6502 program? We define viability as
exponential spread: the program, when placed in cell $(0,0)$ of an
$8 \times 8$ board, fills more than half the cells within 80 scheduler
passes.

The minimal replicator is 8 bytes:
\begin{equation}
\underbrace{\hex{B5}}_{\text{LDA zp,X}}
\underbrace{\hex{00}}_{\text{addr}}
\underbrace{\hex{9D}}_{\text{STA abs,X}}
\underbrace{\hex{00}}_{\text{addr}}
\underbrace{\hex{04}}_{\text{page}}
\underbrace{\hex{E8}}_{\text{INX}}
\underbrace{\hex{90}}_{\text{BCC}}
\underbrace{\hex{F8}}_{\text{offset}}
\end{equation}
with variants at the INX/DEX and branch positions. Longer programs
(9--32 bytes) can include ``slide'' bytes---single-byte opcodes
inserted between core positions---that may or may not preserve viability.

\section{WFST Architecture}

\subsection{Composed Machine}

The scoring transducer is a product-state composition of two reviewers:

\begin{itemize}
\item \textbf{Opcode reviewer} (13 states): Recognizes the structural
  pattern \texttt{LDA~zp,X / STA~abs,X / page / INX|DEX / branch / offset}
  with four slide states allowing variable-length insertions.

\item \textbf{Offset counter} (50 states): Silent byte-position counter.
  After composition, offset verdicts are injected at the product states
  $(\text{seen-branch}, \text{pos}_K)$: \PASS\ if the byte equals
  $(255 - K) \bmod 256$, \FAIL\ otherwise.
\end{itemize}

The composed machine has $13 \times 50 = 650$ states. Additional hard
constraints are injected post-composition:
\begin{itemize}
\item \textbf{Address}: \FAIL\ at states $(\text{seen-LDA}, \text{pos}_K)$
  and $(\text{seen-STA}, \text{pos}_K)$ for any byte $\neq \hex{00}$.
  Simulation revealed that non-zero addresses cause byte~0 to wrap
  into the stack page (8-bit vs.\ 16-bit addressing asymmetry).
\item \textbf{Branch}: \FAIL\ at $(\text{seen-INC}, \text{pos}_K)$
  for branches other than BVC (\hex{50}) and BCC (\hex{90}). Only
  these create infinite loops; BPL/BNE terminate and fail to spread.
\end{itemize}

\subsection{Trainable Weights}

Each non-\FAIL\ transition carries a weight $w_i \in (0, 1]$,
interpreted as $P(\text{this transition is safe})$. At core positions
(LDA, STA, page), the only non-\FAIL\ transition has $w = 1$.
At slide positions, all 256 bytes initially have $w = 1$; training
adjusts these from simulation data.

\subsection{Weighted Forward Algorithm}

The weighted Forward count is:
\begin{equation}
\alpha_s(n) = \sum_{b=0}^{255} w(s, b) \cdot \alpha_{\delta(s,b)}(n-1)
\end{equation}
where the sum excludes \FAIL-verdict transitions, and
$\alpha_s(0) = \mathbf{1}[s \in \mathcal{A}]$ for accept states
$\mathcal{A}$. The WFST's estimate of $P(\text{viable at length } L)$
is $\alpha_{s_0}(L) / 256^L$, giving
$\Beff = -\log_2 \alpha_{s_0}(L) + 8L$.

\section{Training: AND-Gate Blame Assignment}

\subsection{Generative Model}

Each transition $i$ on a path independently passes with probability
$w_i$. The sequence is viable iff \emph{all} transitions pass (AND gate):
\begin{equation}
P(\text{viable} \mid \mathbf{w}) = \prod_{i \in \text{path}} w_i
\end{equation}

\subsection{E-Step: Posterior Blame}

Given that the oracle reports the sequence as non-viable, the posterior
probability that transition $i$ failed is:
\begin{equation}\label{eq:blame}
P(\FAIL_i \mid \text{non-viable}) =
\frac{1 - w_i}{1 - \prod_{j \in \text{path}} w_j}
\end{equation}

When all weights are equal, blame is distributed equally among all
transitions. When one weight is much lower, it receives proportionally
more blame. This is the proper Forward-Backward computation for the
AND-gate topology.

\subsection{M-Step}

For each transition $(s, b)$, accumulate expected counts:
\begin{align}
\hat{p}_{s,b} &= \sum_{\text{viable}} 1
+ \sum_{\text{non-viable}} P(\PASS_i \mid \text{non-viable}) \\
\hat{f}_{s,b} &= \sum_{\text{non-viable}} P(\FAIL_i \mid \text{non-viable})
\end{align}
with Laplace smoothing:
\begin{equation}
w_{s,b}^{\text{new}} = \frac{\hat{p}_{s,b} + \alpha}
{\hat{p}_{s,b} + \hat{f}_{s,b} + 2\alpha}
\end{equation}

\section{Importance Sampling Validation}

The WFST may be miscalibrated. We validate with importance sampling:
sample $x_1, \ldots, x_N$ from the WFST proposal $q$, simulate each,
and compute:
\begin{equation}
\hat{P}(\text{viable}) = \frac{1}{N} \sum_{i=1}^{N}
\mathbf{1}[\text{viable}(x_i)] \cdot \frac{P_{\text{uniform}}(x_i)}{q(x_i)}
\end{equation}
where $q(x) = \prod_{t=1}^{L} \frac{w(s_t, x_t) \cdot \alpha_{\delta(s_t, x_t)}(L-t-1)}{\alpha_{s_t}(L-t)}$ and $P_{\text{uniform}}(x) = 256^{-L}$.

\begin{table}[h]
\centering
\caption{Calibration: WFST vs.\ IS ground truth (after training)}
\begin{tabular}{ccccc}
\toprule
$L$ & WFST $\Beff$ & IS $\Beff$ & Gap (bits) & Viable rate \\
\midrule
8  & 62.0 & 62.0 & 0.0  & 100\% \\
9  & 59.0 & 61.8 & +2.7 & 36\% \\
10 & 57.2 & 61.5 & +4.3 & 12\% \\
11 & 55.9 & 60.2 & +4.3 & 8\% \\
12 & 54.9 & 59.8 & +4.9 & 4\% \\
\bottomrule
\end{tabular}
\label{tab:calibration}
\end{table}

\section{Cumulative $\Beff$ Across Lengths}

A random $L$-byte stream is viable if its first $l$ bytes form a viable
$l$-byte program for any $8 \le l \le L$ (infinite-loop replicators
never execute trailing bytes). Since the offset values at different
lengths are distinct, these events are mutually exclusive:
\begin{equation}
P(\text{viable at length} \le L) = \sum_{l=8}^{L} P_l
\end{equation}

The per-length probability $P_l$ for programs with $k = l - 8$ slide
bytes involves the structural count (from Forward) times the fraction
of safe slide combinations:
\begin{equation}
P_l \approx \frac{Z_{\text{struct}}(l)}{256^l} \cdot p_{\text{safe}}^k
\end{equation}
where $p_{\text{safe}} \approx 26.5/256 \approx 0.104$ (24 safe +
5 risky single-byte opcodes at effective weight 0.5). Since
$p_{\text{safe}} < 1$, the series converges geometrically:
\begin{equation}
\sum_{l=8}^{\infty} P_l = P_8 \cdot \sum_{k=0}^{\infty}
\binom{k+3}{3} \left(\frac{p_{\text{safe}}}{1}\right)^k
\cdot r^k = \frac{P_8}{(1 - p_{\text{safe}} \cdot r)^4}
\end{equation}
where $r$ accounts for the growth of the structural count per
additional slide byte. Numerically, the series converges by
$L \approx 12$, giving cumulative $\Beff \approx 60.9$ bits.

\section{Active Learning}

Simulation is expensive; Forward/sampling is cheap. Active learning
allocates simulation budget to transitions with maximum uncertainty,
measured by the Bernoulli entropy $H(w_i) = -w_i \log_2 w_i -
(1-w_i) \log_2(1-w_i)$.

\begin{algorithm}
\caption{Active Learning Iteration}
\begin{algorithmic}[1]
\State Identify transitions with $H(w_i) > \tau$
\State Temporarily boost their weights by factor $\beta$
\State Sample $N$ sequences from boosted WFST
\State Restore original weights
\State Simulate samples, compute rewards
\State M-step: update weights with AND-gate blame (Eq.~\ref{eq:blame})
\State IS validation: estimate true $\Beff$
\end{algorithmic}
\end{algorithm}

\section{Simulation-Discovered Constraints}

Several constraints were discovered by simulation rather than
theoretical analysis:

\begin{enumerate}
\item \textbf{Address must be \hex{00}}: The LDA zero-page instruction
  uses 8-bit address wrapping, but STA absolute uses 16-bit arithmetic.
  For addr $= A > 0$ and $X = 256 - A$, the source reads byte~0
  (wrapping) but the destination writes to offset $256$ (stack page).
  Byte~0 is physically unable to reach neighbor offset~0.
  Fidelity degrades linearly: $\text{fidelity} = \max(0, (8-A)/8)$.

\item \textbf{Only BVC/BCC spread}: BPL and BNE create finite loops
  (128 and 256 iterations respectively). After the loop, PC falls
  through to byte~8 (typically BRK), ending execution. These copy
  but don't sustain exponential growth (spread $\approx 5$ vs.\ 63).

\item \textbf{Position-dependent slide safety}: The same opcode has
  different viability at different slide positions. TYA (\hex{98})
  is safe as a prefix but lethal between addr and STA (corrupts~A
  before the store). RTS (\hex{60}) is safe as prefix but lethal
  between page and INX.

\item \textbf{Multi-slide interactions}: SEC (\hex{38}) sets carry,
  killing BCC. But SEC followed by CLC (\hex{18}) restores carry=0,
  rescuing BCC. This non-additive interaction requires state in
  the model (or, equivalently, context-dependent weights).
\end{enumerate}

\section{Results}

\begin{table}[h]
\centering
\caption{Effective information content of the minimal replicator}
\begin{tabular}{lcc}
\toprule
Constraint & Bits & Source \\
\midrule
LDA = \hex{B5} & 8.0 & Structural \\
addr = \hex{00} & 8.0 & Simulation \\
STA = \hex{9D} & 8.0 & Structural \\
addr match & 8.0 & Structural \\
page = \hex{04} & 8.0 & Structural \\
inc $\in$ \{INX, DEX\} + branch $\in$ \{BVC, BCC\} & 14.0 & Simulation \\
offset (length-dependent) & 8.0 & Structural \\
\midrule
Total (single length) & 62.0 & \\
Cumulative (all lengths $\le 32$) & $\approx 60.9$ & \\
\bottomrule
\end{tabular}
\end{table}

The IS-validated cumulative $\Beff$ across lengths 8--32 is
approximately \textbf{60.9 bits}, meaning roughly 1 in $2^{61}
\approx 2.3 \times 10^{18}$ random byte streams of length $\ge 8$
encode a viable self-replicating program.

\end{document}
